% Conversão fiel de chapter-05-text-h-0078; origem RELATORIO_ATUAL.odt.
\hypertarget{__RefHeading___Toc4805_3118176633}{}
\refstepcounter{subsection}
\subsection*{4.2.2 Banco de dados utilizado}
\label{ages-iii:text-h-0078}
\addcontentsline{toc}{subsection}{4.2.2 Banco de dados utilizado}

% SOURCE text-p-1039
\originaltab{}O sistema utiliza banco de dados relacional PostgreSQL\space\cite{ages-iii-postgresql}, hospedado em ambiente gerenciado na AWS por meio do Amazon RDS\space\cite{ages-iii-rds}. A modelagem foi estruturada para contemplar as principais entidades do projeto, como usuários, veículos, tags RFID\space\cite{ages-iii-rfid-nist}, acessos e transações, garantindo integridade relacional e consistência dos dados.\par

% SOURCE text-p-1040
Além do armazenamento operacional, a estrutura permite consultas voltadas ao dashboard administrativo, viabilizando métricas de fluxo de veículos, histórico de acessos, ocupação e demais indicadores relevantes para a gestão do estacionamento.\par

% FIGURE draw-image-0017
\memorialfigure{imagens/original/asset-0014.png}{\textbf{Figura 14 –} Diagrama Banco de Dados - AGES III}{ages-iii:draw-image-0017}{Fonte: pendente de identificação}{0.92} % <<ASSET draw-image-0017 asset-0014>>
\indent Mais detalhes podem ser encontrados aqui: \href{https://tools.ages.pucrs.br/sistema-de-controle-inteligente-de-estacionamento-com-rfid/impinj-rfid-wiki/-/wikis/Banco\%20de\%20Dados}{Banco de Dados}\par
